\section{开放挑战与研究路线}

\subsection{课程给出的长期挑战}
课程后段将挑战集中在三个层面：认知层（更可靠的推理）、系统层（更稳健的记忆与规划）、组织层（更可信的评测与部署）。

\begin{figure}[H]
\centering
\includegraphics[width=0.84\textwidth]{slide-06-open-challenges.jpg}
\caption{视频约 01:24:30 讲义页：未来方向聚焦在可靠推理、记忆治理与长程规划}
\end{figure}

\begin{importantbox}{未来两年的关键技术杠杆}
\begin{enumerate}
    \item 可验证推理：把 reasoning 结果与形式化检查器连接；
    \item 记忆治理：解决污染传播、版本冲突和过期淘汰；
    \item 规划泛化：提升跨任务迁移能力，降低每任务重调成本。
\end{enumerate}
\end{importantbox}

\begin{knowledgebox}{为何“Reasoning + Memory + Planning”必须联合优化}
单点优化很容易出现木桶效应：推理增强但记忆劣化，会放大错误；规划增强但验证器弱，会更快走偏。联合优化的目标是系统级稳定收益，而非局部峰值。
\end{knowledgebox}

\begin{warningbox}{部署中的现实约束}
企业场景往往存在严格延迟预算、权限边界和审计要求。研究原型若忽略这些约束，迁移到生产时会出现显著性能回撤。
\end{warningbox}

\subsection{本章小结}
课程给出的路线不是追逐单一 SOTA，而是建设可验证、可回滚、可扩展的 Agent 系统栈。长期竞争力来自系统工程能力。

